\documentclass[11pt]{extarticle}
\setcounter{secnumdepth}{4}
\usepackage{times}
\setcounter{tocdepth}{4}
\usepackage[english]{babel}
\usepackage{float}
\usepackage{amsmath}
\usepackage{subcaption}
\usepackage{algorithm}
\usepackage{algpseudocode}
\usepackage{pdflscape}
\usepackage{titlesec}
\usepackage{textcomp}
\usepackage{tikz}
\usetikzlibrary{calc, positioning, arrows.meta, shapes.multipart}
\usepackage{changepage}
\usepackage{multirow}
\usepackage[a4paper, margin=2.5cm]{geometry}
\usepackage{titlepic}
\usepackage{multicol}
\usepackage{makecell}
\makeatletter
\usepackage{graphicx}
\usepackage{authblk}
\usepackage{blkarray,booktabs,bigstrut}
\usepackage[colorlinks=true, allcolors=blue]{hyperref}
\usepackage{titling}
\usepackage[linewidth=1pt]{mdframed}
\usepackage{lipsum}
\usepackage{outlines}
\usepackage{tocloft}
\usepackage{caption}

%Coloring for table:
\usepackage{colortbl}
\usepackage{array}
\definecolor{badred}{HTML}{F4CCCC}
\definecolor{goodgreen}{HTML}{D9EAD3}
\definecolor{bestgreen}{HTML}{93C47D}

% Short macros to keep the table code clean
\newcommand{\cR}[1]{\cellcolor{badred}#1}
\newcommand{\cG}[1]{\cellcolor{goodgreen}#1}
\newcommand{\cB}[1]{\cellcolor{bestgreen}\textbf{#1}} % Dark green + Bold for >40%

\renewcommand\maketitlehooka{\null\mbox{}\vfill}
\renewcommand\maketitlehookd{\vfill\null}

\titleformat{name=\section,numberless}[display]
  {\normalfont\huge\bfseries\filcenter}
  {}
  {0pt}
  {\centerline{\rule{\textwidth}{0.5pt}}\vspace{1.5ex}}
  [\vspace{0.5ex}\centerline{\rule{\textwidth}{0.5pt}}]

\titlespacing*{\section}{0pt}{15pt}{10pt}

\titleformat{\paragraph}
{\normalfont\normalsize\bfseries}{\theparagraph}{1em}{}

\usepackage{fancyhdr}
\pagestyle{fancy}
\fancyhf{} % Clear all header and footer fields
\fancyfoot[C]{\thepage} % Set page number to the right in the footer
\renewcommand{\headrulewidth}{0pt} % Optional: removes the header line

\begin{document}

% =========== TITLE PAGE =================

\begin{titlepage}
    \begin{center}
        \begin{tikzpicture}[remember picture, overlay]
            \draw[line width=1pt]
            ($(current page.north west) + (2cm,-2cm)$)
            rectangle
            ($(current page.south east) + (-2cm,2cm)$);
        \end{tikzpicture}

        {\large \textbf{LA ROCHELLE UNIVERSITÉ}} \\
        \vspace{0.5cm}
        {\large Licence Informatique} \\
        \vspace{1cm}
        
\includegraphics[width=0.4\textwidth]{figures/La_Rochelle_Universite.png}
        \vspace{1cm}

        {\fontsize{27}{33}\selectfont \textbf{INTERNSHIP REPORT}}

        \vspace{1cm}
        {\Large \textbf{TRAN Minh Duong\\}}

        \vspace{0.8cm}
        {\large \textit{Supervised by:}\\}
        \vspace{0.2cm}
        {\large \textbf{Dr. Emanuela BOROS} \\
            \textbf{Prof. Antoine DOUCET} \\
            \normalsize(Laboratoire L3i)\\}


        \vspace{1.5cm}
        {\Large Title:\\}
        \vspace{0.5cm}
        {\fontsize{20}{25}\selectfont \textbf{``Historical Fidelity in OCR Post-Correction: A Comparative Study of Fine-Tuned LLMs and Seq2Seq Models''}}\\[2cm]

        \vspace{1cm}
        {\Large Academic Year: 2025-2026}
    \end{center}
    \vspace*{\fill}
\end{titlepage}

\pagestyle{fancy} % Enable numbering from this point forward
\pagenumbering{arabic}

\tableofcontents
\newpage

\setlength{\parindent}{0pt}
\setlength{\parskip}{1em}
\newpage


% ============ ACKNOWLEDGEMENTS ============
\section*{ACKNOWLEDGEMENT}
% \addcontentsline{toc}{section}{Acknowledgement}

I would like to express my deepest gratitude to all those who have mentored and supported me throughout this internship.

I want to sincerely thank Dr. Emanuela BOROS for her patient guidance for the past few months. Had it not been for her encouragement and advice, I could not have completed my thesis. Furthermore, I would like to convey my deepest appreciation to Prof. Antoine DOUCET for providing me with the invaluable opportunity to intern at La Rochelle University's L3i laboratory. I would be negligent if I did not mention all the people in the lab for their warmest welcome and hospitality.

I am also extremely grateful to take part in the double degree exchange program between La Rochelle University and the University of Science and Technology of Hanoi. I would like to express my gratitude to all of the lecturers at La Rochelle University for their support and guidance despite my limited fluency in French. This experience has allowed me to connect with people from different cultural backgrounds and reshaped the way I understand and engage with the world around me.

To my lovely family, thank you for all your blessings, I love you all. To my dear Oanh, thank you for the enormous emotional support. Thank you Le Linh Long for generously sponsoring my Overleaf pro plan and guiding me with my writings.

Last but not least, to those who are willingly spending your time reading my work, thank you and I hope you will have a wonderful time as I had when writing it.

\newpage

% ===================================== ABSTRACT ==============================================
\section*{ABSTRACT}
\addcontentsline{toc}{section}{Abstract}

Digitizing historical documents is often difficult due to significant character and word errors introduced during the Optical Character Recognition (OCR) process. While modern AI models offer powerful correction capabilities, they often struggle to distinguish between actual OCR noise and valid historical orthography, leading to the risk of "over-correction."

This study evaluates two different architectural paradigms to solve this problem: a standard encoder-decoder model (BART) and a modern Large Language Model (Qwen) trained efficiently using QLoRA. We further investigate whether enriching model inputs with document-level metadata (publication year, language, and collection origin) can assist model predictions with historical context and reduce over-correction. 
Using the multilingual HIPE-OCRepair 2026 benchmark, we analyze the trade-off between error reduction and historical fidelity.

Our results show that while Large Language Models (LLMs) achieve better semantic understanding, they are more prone to aggressive paraphrasing and hallucinations. Metadata conditioning substantially mitigates these generative biases in smaller models, though its effect on larger models remains inconsistent. In contrast, standard denoising encoder-decoder models (Seq2Seq) maintain tighter lexical constraints at the cost of contextual fluidity. Finally, we propose a qualitative error taxonomy to categorize these behaviors, helping researchers understand how to safely use generative AI to clean up historical archives without losing their original character.

\textbf{Keywords}: OCR Post-Correction, Historical Documents, Large Language Models (LLMs), Sequence-to-Sequence (Seq2Seq), Parameter-Efficient Fine-Tuning (PEFT), Over-correction

\newpage

\section{Introduction}
\subsection{Presentation of the Host Institution}
Established in 1993, the \textbf{Laboratoire Informatique, Image, Interaction (L3i)} is the primary research center for Digital Sciences at La Rochelle Université \cite{l3i_website}. The laboratory consists of approximately 98 members specializing in fields such as Computer Science, Image Processing, and Human-Computer Interaction. The L3i is known for its transversal research vision, combining expertise in Artificial Intelligence and Knowledge Management with other disciplines, including Environmental Sciences and the Humanities.

The laboratory maintains a dynamic ecosystem of industrial and academic partnerships. It manages several joint laboratories, such as IDEAS (with Yooz), which focuses on AI applied to administrative documents, and SAIL (avec Mangas.IO), dedicated to digital comics analysis. On an international level, the L3i collaborates closely with various institutions, including the ICTLab (USTH) in Vietnam. This connection is particularly significant for this research, as it strengthens scientific exchanges between Europe and Southeast Asia.

Within this research environment, my internship is supervised by Dr. Emanuela Boros and Prof. Antoine Doucet. The work is deeply integrated into the framework of HIPE-OCRepair\cite{hipe_ocrepair_2026}, an upcoming competition for the ICDAR 2026 conference \cite{icdar_2026}. This initiative is dedicated to the OCR post-correction of multilingual historical documents using Large Language Models (LLMs). The project aims to address the "OCR debt" found in large digitized collections, where manual correction is not possible. By providing standardized datasets in English, French, and German, the HIPE-OCRepair benchmark allows for a systematic evaluation of different AI architectures for preserving historical text and orthography.

\subsection{Context and Problem Statement}
As the world moves toward digital records, historical documents are also being digitized. These archives are important sources for researchers to understand past events, social changes, and cultural perspectives. While scanning and Optical Character Recognition (OCR) technology improve the accessibility of these documents, the extracted text is often noisy and full of errors. The quality of OCR frequently suffers due to physical degradation over time (such as bleed-through, ink spills, and fading), poor print quality, outdated typefaces, and complex page layouts. Consequently, this noise significantly hinders text mining techniques and keyword searches, making it difficult for scholars to extract meaningful information. To unlock the full potential of these primary sources, post-OCR correction techniques, such as refining and enhancing the generated text, are crucial. 

However, working with historical text involves more than just reading degraded scans. Historical documents frequently contain archaic spellings, unique grammar, and vocabulary that are no longer used today. In recent years, Large Language Models (LLMs) have significantly advanced the field of natural language processing, demonstrating remarkable success in analyzing complex context and generating coherent text with minimal training \cite{liang2023holisticevaluationlanguagemodels, zhao2026surveylargelanguagemodels}. Yet, applying these state-of-the-art models to historical OCR post-correction introduces significant challenges. Standard language models are primarily trained on modern text. Without specialized training, they tend to "over-correct" the text, replacing valid historical words with modern equivalents. Furthermore, distinguishing between actual OCR noise and a valid historical spelling variation is a highly complex task. When faced with degraded text, models can struggle to guess the correct word, sometimes generating text that alters the original meaning. Therefore, there is a critical need to evaluate how well different model architectures can reduce scanning errors while keeping the historical authenticity of the document intact.

\subsection{Proposed Solution and Research Question}

To address the dual challenge of reducing OCR noise and preventing over-correction, this thesis evaluates several modern deep learning models. Instead of relying on a single architecture, this study compares different approaches across four main experiments:

\begin{itemize}
    \item \textbf{Sequence-to-Sequence (Seq2Seq) Models:} We evaluate the standard encoder-decoder approach by fine-tuning models from the BART family \cite{lewis2019bartdenoisingsequencetosequencepretraining}. This includes testing different model sizes (Base vs. Large) and exploring multilingual models (mBART).
    
    \item \textbf{Decoder-Only (Generative) Models:} We investigate modern Large Language Models (LLMs), such as Qwen3 \cite{yang2025qwen3technicalreport}. We compare different ways to use them, from simple Zero-Shot prompting to Parameter-Efficient Fine-Tuning (using QLoRA).
    
    \item \textbf{Domain-Specific Models:} We compare our models against systems that were specifically designed for OCR correction. This includes specialized Seq2Seq models (like Pykale, the correction model released by Thomas et al. \cite{thomas-etal-2024-leveraging}) and specialized generative models (like PleIAs \cite{pleIAs}).
    
    \item \textbf{The Impact of Metadata:} We conduct specific experiments to see if adding explicit historical context (like publication year, language, and dataset origin) to the prompts improves the models' ability to correct text accurately.
\end{itemize}

By comparing these different methods through both quantitative metrics and qualitative analysis, this study explores the current state of OCR post-correction. Ultimately, all of these experiments are designed to answer one central, falsifiable hypothesis:

\begin{mdframed}[backgroundcolor=gray!10, linewidth=0pt, innertopmargin=10pt, innerbottommargin=10pt, innerleftmargin=10pt, innerrightmargin=10pt]
\textbf{Central Research Question:} \textit{Does parameter-efficient fine-tuning on domain-specific historical data allow modern LLMs to reduce OCR noise while preserving archaic orthography significantly better than traditional, fully fine-tuned Seq2Seq models?}
\end{mdframed}


%================================= BACKGROUND & RELATED WORK ======================================
\section{Background and Related Work}

\subsection{OCR in Historical Documents}
Optical Character Recognition (OCR) is a method of extracting textual content from images of documents. This procedure is frequently required when a digital copy of the text is unavailable, such as when processing physical archives. For decades, massive volumes of scanned documents have been processed using traditional OCR algorithms, driving major digitization efforts by Google Books \cite{leetaru_2008}, the Europeana Newspapers Project \cite{neudecker_2016}, and the National Library of Australia \cite{holley_2009}. 

While modern state-of-the-art models based on deep learning have significantly improved baseline accuracy, substantial noise often persists in the output. This residual noise severely degrades the performance of downstream natural language processing tasks, such as named entity recognition, part-of-speech tagging, and text summarization \cite{10.1007/978-3-030-64452-9_3,vanStrien2020AssessingTI}. These limitations are particularly pronounced when OCR is applied to historical documents. Unlike modern text, historical archives present a unique set of interconnected challenges that heavily disrupt OCR engines:
\begin{itemize}
    \item \textbf{Physical Degradation:} Centuries-old documents frequently suffer from severe physical aging, including faded ink, paper discoloration, and ink bleed-through from the opposite side of the page.
    \item \textbf{Archaic Typography:} Historical texts often feature non-standard or obsolete fonts (such as Fraktur in German documents), as well as deteriorated lead types and inconsistent spacing.
    \item \textbf{Layout Complexity:} The structure of historical newspapers and books often includes dense, multi-column layouts, marginalia, and inconsistent reading orders that confuse standard OCR bounding boxes.
\end{itemize}

Because of these physical and structural ambiguities, the resulting OCR output is often so degraded that even human interpretation becomes subjective. Consequently, robust post-correction pipelines are required to reconstruct the text before it can be effectively analyzed by scholars.

\subsection{OCR Post-Correction as a Task}
To mitigate the limitations of raw OCR engines, early post-correction techniques primarily relied on rule-based methods and dictionary lookups to fix out-of-vocabulary words. However, these traditional systems struggle significantly with historical documents. Because historical texts are full of archaic spellings, static dictionaries often flag valid historical words as errors, unintentionally replacing them with modern equivalents.

To address these limitations, the field moved toward data-driven machine learning, reframing OCR correction as a Machine Translation (MT) problem. A key driver of this shift was the ICDAR Post-OCR Text Correction competitions (Chiron et al., 2017 \cite{chironicdar2017}; Rigaud et al., 2019 \cite{rigaudicdar2019}), which exposed the shortcomings of statistical methods and promoted the adoption of Neural Machine Translation (NMT). 

In this paradigm, neural networks map noisy input text directly to clean output text. Dong and Smith (2018) \cite{dong-smith-2018-multi} showed that character-level sequence-to-sequence models with attention are effective for unsupervised OCR correction. Schaefer and Neudecker (2020) \cite{schaefer-neudecker-2020-two} extended this with a two-step NMT method: separating detection and correction phases for historical printed documents. By casting correction as sequence-to-sequence translation, these models move beyond isolated word lookups and learn complex, context-dependent error patterns caused by historical document degradation. 

Recent studies have further expanded this sequence-to-sequence formulation to highly diverse document types and resource conditions. Post-correction has recently been explored for early modern theatre \cite{debaene-etal-2025-evaluating} and endangered Indigenous languages of Latin America \cite{agarwal2025ailla}. Furthermore, because aligned OCR-to-ground-truth pairs are historically scarce, synthetic data generation has emerged as a crucial strategy for pre-training these post-correction models \cite{guan2024advancing,bourne2025scrambled}. These advancements prove that post-correction remains highly dependent on language, domain, and the availability of specialized training data.

\subsection{Neural Approaches for Text Correction}
To address this translation task, modern research relies on deep learning architectures. This study focuses on comparing the two dominant paradigms in the field.

\subsubsection{Sequence-to-Sequence Models (BART)}
Sequence-to-Sequence (Seq2Seq) architectures, such as BART (Lewis et al. \cite{lewis2019bartdenoisingsequencetosequencepretraining}), have become a standard baseline for text correction. Utilizing a standard transformer-based structure, BART operates as a denoising encoder-decoder. It is pre-trained by corrupting text with an arbitrary noising function and learning to reconstruct the original text. 

The encoder processes the entire noisy OCR sequence to build a bidirectional context, while the autoregressive decoder generates the clean output. Unlike BERT, BART's autoregressive decoder allows it to generate outputs of varying lengths, making it highly effective at handling OCR insertions and deletions. The efficacy of this approach was notably demonstrated by Soper et al. (2021) \cite{soper-etal-2021-bart}, who were the first to successfully utilize sentence-level transformer models (specifically fine-tuned BART) for the post-correction of OCR newspaper text, framing the problem directly as a translation task. However, despite this flexibility, it still functions as a highly constrained local translator, anchoring its generation strictly to the underlying input sequence rather than engaging in free-form generation.

Furthermore, sequence-to-sequence models are not inherently preservation-oriented. Their pre-training data and subword tokenization strategies heavily bias them toward frequent modern forms, especially since rare historical spellings were largely underrepresented in their training corpora \cite{soper-etal-2021-bart}. As a result, a model may successfully reduce word-level errors while simultaneously increasing character-level divergence from the ground truth through aggressive modernization and lexical substitution.

\subsubsection{Large Language Models (LLMs) and Hallucinations}
More recently, autoregressive Large Language Models (LLMs) with decoder-only architectures, such as Qwen3 (introduced by Alibaba in 2025 \cite{yang2025qwen3technicalreport}), have been applied to post-correction, with researchers actively exploring their potential to unlock historical newspaper archives (Thomas et al., 2024 \cite{thomas-etal-2024-leveraging}). These models possess a massive capacity to understand broad semantic context, allowing them to guess missing words in highly degraded text.

However, because LLMs are primarily pre-trained on modern internet corpora, their generative nature poses significant risks for historical documents. They frequently suffer from "over-correction," in which the model automatically modernizes valid historical spellings to conform to up-to-date grammar rules. They are also highly prone to "hallucinations." \textbf{In the context of historical OCR post-correction, a hallucination is the generation of fluent, historically plausible, but optically unsupported text.} This token- or sentence-level invention occurs when the model prioritizes its internal semantic probability (what "makes sense") over the visual constraints of the original noisy string (what was actually scanned). This semantic drift creates a fundamental trade-off between achieving reading fluency and preserving archival authenticity.

\subsubsection{Parameter-Efficient Fine-Tuning (PEFT)}
As Large Language Models (LLMs) emerged as dominant NLP tools, researchers began evaluating their out-of-the-box capabilities for text restoration. However, as demonstrated by Boros et al. (2024) \cite{boros-etal-2024-post} in an extensive exploratory study of fourteen foundation models, standard zero-shot and few-shot prompting techniques are largely insufficient for this task. Their quantitative and qualitative analyses revealed that without domain-specific fine-tuning, general-purpose LLMs struggle to balance OCR noise reduction with the preservation of historical nuances.

To overcome these limitations and prevent LLMs from over-correcting, several mitigation strategies have recently been proposed. These include training on synthetic data (Guan et al. 2024 \cite{guan2024advancing} and Bourne et al. 2025 \cite{bourne2025scrambled}), utilizing reference-based correction (Do et al. \cite{do2025reference}), applying constrained decoding (Sastre et al. \cite{sastre2025postocr}), and incorporating multimodal evidence directly from document images (Grieif et al. \cite{greif2025multimodal}). However, for text-only workflows, the models must be specialized directly on historical corpora. Traditionally, full fine-tuning requires massive computational resources that are unavailable to standard research laboratories. Parameter-Efficient Fine-Tuning (PEFT) techniques, specifically QLoRA (Quantized Low-Rank Adaptation) \cite{dettmers2023qloraefficientfinetuningquantized}, solve this hardware bottleneck. By freezing the original model weights and training only a small, highly compressed set of new parameters, QLoRA allows massive models to adapt to specific tasks using a single standard GPU.

\subsection{Metadata in Document Analysis}
Historical documents are rarely context-free. Their accurate interpretation depends heavily on variables such as language, publication date, document type, and collection origin \cite{vanstrien2020assessing}. Such metadata provides vital signals about expected spelling conventions, named entities, vocabulary, and even the specific OCR noise patterns generated by different digitization workflows. 

In the context of OCR post-correction, metadata serves as a theoretical anchor to help models distinguish between true OCR degradation and valid historical variation. For example, a document's publication date can explicitly indicate whether an archaic spelling is temporally plausible, while the language tag constrains character patterns and morphological rules. Despite its wide availability in large digital libraries, metadata-aware conditioning for neural OCR post-correction remains largely underexplored. This thesis directly addresses this gap by evaluating whether providing these explicit contextual constraints can prevent neural architectures from over-correcting historical texts.

%========================================== METHODOLOGY =================================================
\section{Methodology and Experimental Setup}

\subsection{Task Definition}
The primary objective of this study is to perform OCR post-correction on degraded historical documents. Formally, this is defined as a sequence mapping task where a noisy OCR input sequence $X = (x_1, \dots, x_n)$ must be transformed into a corrected target sequence $Y = (y_1, \dots, y_m)$ that matches the curated transcription.

In this study, we evaluate the models under two distinct conditioning settings. In the standard \textbf{metadata-free} setting, the model relies solely on the corrupted text to predict the clean sequence, defined as:
\begin{equation}
    \hat{Y} = f_{\theta}(X)
\end{equation}
where $f_{\theta}$ is the neural model mapping the noisy text to a corrected output.

To test the impact of document-level context, we also define a \textbf{metadata-aware} OCR post-correction setting:
\begin{equation}
    \hat{Y} = f_{\theta}(X, M)
\end{equation}
where $M$ represents available historical annotations, such as publication year, language, and dataset origin. 


\subsection{HIPE-OCRepair Benchmark}
To train and evaluate the models, this study utilizes the HIPE-OCRepair benchmark \cite{hipe_ocrepair_2026}. This dataset is specifically designed for the ICDAR 2026 competition and provides a standardized environment for OCR post-correction on historical documents. It covers three languages: English, French, and German, originating from established archives such as DTA-19, Impresso, Overproof, and subsets of previous ICDAR 2017 competitions.

\subsubsection{Data Format and Structure}
A key characteristic of this dataset is that the models do not have access to the original scanned images; they must perform the correction purely based on text and context. The data is provided in JSONL format, where each text segment contains the \textbf{OCR Hypothesis} (the noisy input), the \textbf{Curated Ground Truth} (the clean target), and \textbf{Document Metadata} (language, publication date, and origin).

\subsubsection{Correction Challenges and Baseline Statistics}
Historical OCR noise extends beyond simple misspellings, encompassing arbitrary spaces, incorrect capitalization, and character replacements caused by physical degradation. Table \ref{tab:dataset_examples} illustrates these challenges, highlighting the differences between the noisy OCR input and the clean ground truth. 

To understand the initial degradation levels, we conducted an Exploratory Data Analysis (EDA) across the data splits (see Table \ref{tab:eda_summary}). Baseline Character Error Rate (CER) and Word Error Rate (WER) vary significantly depending on the historical collection. 

\begin{table}[H]
    \centering
    \renewcommand{\arraystretch}{1.5}
    \begin{tabular}{>{\raggedright\arraybackslash}p{2.5cm} >{\raggedright\arraybackslash}p{6cm} >{\raggedright\arraybackslash}p{6cm}}
    \toprule
    \textbf{Subset} & \textbf{OCR Transcription (Input)} & \textbf{Ground Truth (Target)} \\
    \midrule
    
    icdar2017 (FR) & 
    Blondel, in\textbf{apertc}ur général dès finances... & 
    Blondel, in\textbf{specte}ur général dès finances... \\
    
    icdar2017 (FR) & 
    on du maréchal Pétain dén\textbf{tant 4 V  }ex-général Giraud tout... & 
    on du maréchal Pétain dén\textbf{iant, a l'}ex-général Giraud tout... \\
    
    impresso-snippets (FR) & 
    Suv ce nombre, un peu inférieur \textbf{ù} celui du mois \textbf{précèdent}, il y a eu 7 & 
    Suv ce nombre, un peu inférieur \textbf{à} celui du mois \textbf{précédent}, il y a eu 7 \\
    
    impresso-snippets (FR) & 
    \textbf{s'ef-.forçait} de calmer les mutins\dots & 
    \textbf{s'efforçait} de calmer les mutins\dots \\
    
    \midrule
    overproof (EN) & 
    A publ\textbf{io mooting} of \textbf{tho rosidonts} and \textbf{ratopaors...} & 
    A publ\textbf{ic meeting} of \textbf{the residents} and \textbf{ratepayers}...\\
    
    icdar2017 (EN) & 
    t\textbf{lie} person who witnessed \textbf{tl.e} fact had not an opportunity of...& 
    t\textbf{he} person who witnessed \textbf{the} fact had not an opportunity of...\\
    
    \midrule
    dta19 (DE) & 
    ...ettergange neben ihrem Kü\textbf{dm}en- und Kupfergeschirr kn... &
    ...ettergange neben ihrem Kü\textbf{ch}en- und Kupfergeschirr kn... \\
    
    impresso-nzz (DE) & 
    ...Sonntag  E\textbf{mm-drerssrgster Iahrga»-  }. Druck und Verlag von O...& 
    ...Sonntag  E\textbf{inunddreissigster Jahrgang}. Druck und Verlag von O... \\
    
    \bottomrule
    \end{tabular}
    \caption{Examples of OCR errors (words in bold) and their target corrections for the different subsets of the HIPE-OCRepair dataset.}
    \label{tab:dataset_examples}
\end{table}

\begin{table}[H]
    \centering
    \renewcommand{\arraystretch}{1.2}
    \resizebox{\textwidth}{!}{
    \begin{tabular}{llcccc}
    \toprule
    \textbf{Split \& Dataset} & \textbf{Period} & \textbf{Languages (Samples)} & \textbf{CER (mean $\pm$ std)} & \textbf{WER (mean $\pm$ std)}\\
    \midrule
    \multicolumn{5}{l}{\textbf{Train Split}} \\
    \midrule
    dta19 & 1815 - 1891 & de (570) & 0.0336 $\pm$ 0.0268 & 0.2369 $\pm$ 0.1827  \\
    icdar2017 & N/A & en (455), fr (391) & 0.0417 $\pm$ 0.0316 & 0.1740 $\pm$ 0.1139 \\
    impresso-nzz & 1781 - 1946 & de (150) & 0.0617 $\pm$ 0.1018 & 0.2477 $\pm$ 0.2789 \\
    impresso-snippets & 1801 - 1959 & en (50), de (50), fr (50) & 0.0367 $\pm$ 0.0297 & 0.2252 $\pm$ 0.1997 \\
    overproof-combined & 1843 - 1954 & en (146) & 0.0866 $\pm$ 0.0821 & 0.3112 $\pm$ 0.1929 \\
    
    \midrule
    \multicolumn{5}{l}{\textbf{Validation (Dev) Split}} \\
    \midrule
    dta19 & 1797 - 1898 & de (330) & 0.0342 $\pm$ 0.0270 & N/A \\
    icdar2017 & N/A & en (188) & 0.0511 $\pm$ 0.0340 & 0.2168 $\pm$ 0.1256 \\
    impresso-snippets & 1801 - 1956 & en (10), de (10), fr (10) & 0.0414 $\pm$ 0.0345 & 0.2561 $\pm$ 0.1892 \\
    overproof-combined & 1872 - 1954 & en (30) & 0.0845 $\pm$ 0.0577 & 0.3427 $\pm$ 0.1666 \\
    
    \midrule
    \multicolumn{5}{l}{\textbf{Test Split}} \\
    \midrule
    dta19 & 1802 - 1863 & de (240) & 0.0342 $\pm$ 0.0270 & 0.1870 $\pm$ 0.1391 \\
    icdar2017 & N/A & en (101), fr (100) & 0.0305 $\pm$ 0.0287 & 0.1330 $\pm$ 0.0918 \\
    impresso-nzz & 1780 - 1939 & de (17) & 0.0586 $\pm$ 0.0870 & 0.2305 $\pm$ 0.3140 \\
    impresso-snippets & 1801 - 1959 & en (100), de (100), fr (100) & 0.0365 $\pm$ 0.0300 & 0.2257 $\pm$ 0.2258 \\
    overproof-combined & 1859 - 1949 & en (32) & 0.1037 $\pm$ 0.1384 & 0.3041 $\pm$ 0.1932 \\
    \bottomrule
    \end{tabular}
    }
    \caption{Exploratory Data Analysis showing the baseline error rates and sample distributions across the different datasets and splits.}
    \label{tab:eda_summary}
\end{table}

\subsubsection{Data Preprocessing and Chunking Strategy}
A major technical challenge in OCR post-correction is sequence length. Historical documents frequently exceed the strict 1024-token context window limitations of transformer models like BART. To address this without losing contextual integrity, we developed a custom preprocessing pipeline to dynamically partition the documents.

Instead of employing naive, rigid truncation, which risks cutting off words mid-sentence and destroying semantic context, the pipeline utilizes a sequence-matching algorithm to identify exact textual anchor points between the noisy OCR sequence and the ground truth. The algorithm partitions the text into manageable chunks (targeting approximately 300 words) while explicitly prioritizing punctuation boundaries (e.g., periods, exclamation marks) within a flexible search window, see Figure \ref{fig:chunking} below. 

\begin{figure}[h]
    \centering
    \includegraphics[width=\textwidth]{report/figures/chunk-pipeline.drawio.pdf}
    \caption{Chunking Strategy based on word anchors and punctuation boundaries.}
    \label{fig:chunking}
\end{figure}

This guarantees that both the input text and the target text remain perfectly aligned and grammatically whole. During the inference phase, these aligned chunks are processed independently by the models and subsequently stitched back together to reconstruct the full document for evaluation. Finally, the partitioned data was isolated into distinct sub-datasets based on language (English, French, German) and aggregated into a combined multilingual dataset to facilitate targeted training.

\subsection{Experimental Models}
We evaluate how different architectures balance error reduction and historical preservation by comparing two neural paradigms. BART is chosen as the Sequence-to-Sequence baseline, as it is the standard in the ICDAR text-correction community. Qwen3 is chosen for the generative paradigm over alternatives like LLaMA3 or Mistral due to its strong multilingual capabilities and instruction-tuning \cite{yang2025qwen3technicalreport}, making it ideal for processing our French and German subsets.

\subsubsection{Seq2Seq Baseline: BART Fine-Tuning}
We employ the BART (Bidirectional and Auto-Regressive Transformers) family of models as our baseline. Because BART utilizes a full encoder-decoder architecture, it is naturally suited for text restoration tasks. The encoder processes the entire noisy OCR sequence to build a bidirectional context, while the decoder autoregressively generates the corrected text. 

In this experiment, BART is fine-tuned strictly as a translation model. It receives no explicit instructions or metadata; it simply learns the statistical mapping from the noisy input domain to the clean target domain. This architecture generally tends to produce outputs closely constrained by the input token sequence due to its denoising pre-training objective.

\subsubsection{LLM Challenger: Qwen Fine-Tuning with QLoRA}
To represent the modern generative approach, we utilize Qwen, an autoregressive decoder-only Large Language Model. Unlike BART, Qwen is instruction-tuned. The model is guided by a system prompt that explicitly defines its persona and task (e.g., \textit{"You are an expert archivist correcting OCR errors in historical documents..."}), followed by the noisy text.

Because fully fine-tuning an LLM with billions of parameters is computationally prohibitive and prone to catastrophic forgetting, we adapt Qwen using Parameter-Efficient Fine-Tuning (PEFT). Specifically, we implement Quantized Low-Rank Adaptation (QLoRA). In this setup, the base weights of the Qwen model are frozen and quantized to 4-bit precision to fit within standard GPU memory limits. We then introduce and train small, low-rank matrices (LoRA adapters) specifically on the HIPE-OCRepair dataset. This allows the LLM to learn the specific orthographic rules of historical texts without losing its massive pre-trained semantic knowledge.

\subsection{Model Training and Setup}
To guarantee reproducibility and an equitable comparison, both models were trained under controlled hardware conditions on L3i’s calculmaster SLURM nodes.

\subsubsection{Hyperparameter Configurations}
Table \ref{tab:hyperparameters} outlines the core hyperparameters utilized during the fine-tuning process. For the BART family, training stabilized over 15 epochs. Numerical instability, such as exploding gradients in \texttt{bart-large}, was successfully resolved by intercepting the backward hook and forcing gradient tensors into \texttt{float32} precision, while maintaining \texttt{fp16} for the forward passes.

For the Qwen generative models, the model weights were quantized to 4-bit NormalFloat (\texttt{nf4}) precision using the \texttt{BitsAndBytes} library. LoRA matrices with a high rank of $r=64$ ($\alpha=16$) were used for all linear target modules. This rank was chosen after initial tests showed underfitting at standard dimensions ($r=16$ and $r=32$) in the noisy historical domain. No further studies on rank were conducted due to hardware constraints.

\begin{table}[H]
    \centering
    
    \begin{subtable}{0.48\textwidth}
        \centering
        \renewcommand{\arraystretch}{1.2}
        \resizebox{\textwidth}{!}{
        \begin{tabular}{lcc}
        \toprule
        \textbf{Parameter} & \textbf{BART-Base} & \textbf{BART-Large} \\
        \midrule
        Epochs & 15 & 15 \\
        Learning Rate & 5e-5 & 2e-5 \\
        Train Batch Size & 2 & 1 \\
        Grad. Accumulation & 8 & 16 \\
        Weight Decay & 0.02 & 0.02 \\
        Warmup Steps & 100 & 100 \\
        Precision & FP16 & FP16 \\
        Max Grad Norm & 1.0 & 1.0 \\
        GPU & RTX 2080Ti & RTX 2080Ti\\
        \bottomrule
        \end{tabular}
        }
        \caption{BART Training Setup}
    \end{subtable}\hfill
    \begin{subtable}{0.48\textwidth}
        \centering
        \renewcommand{\arraystretch}{1.2}
        \resizebox{\textwidth}{!}{
        \begin{tabular}{lc}
        \toprule
        \textbf{Parameter} & \textbf{Qwen3 (4B \& 8B)} \\
        \midrule
        Epochs & 5 \\
        Learning Rate & 2e-4 \\
        Train Batch Size & 8 \\
        Grad. Accumulation & 2 \\
        Optimizer & Paged AdamW (32-bit) \\
        Precision & 4-bit (NF4) + BF16 \\
        LoRA Rank ($r$) & 64 ($\alpha=16$) \\
        LoRA Target Modules & all-linear \\
        GPU & NVIDIA A40, H100 \\
        \bottomrule
        \end{tabular}
        }
        \caption{Qwen QLoRA Setup}
    \end{subtable}
    \caption{Core hyperparameters utilized during the fine-tuning process.}
    \label{tab:hyperparameters}
\end{table}

\subsubsection{Prompt Template and Metadata Integration}
A major hypothesis of this study is that generative models and Seq2Seq architectures can benefit from explicit context. To test this, we developed a secondary data pipeline that injects document metadata (Language, Dataset origin, and Publication Year) directly into the model inputs.

For \textbf{BART}, which operates as a direct text-to-text translator without a conversational template, metadata was added to the beginning of the noisy OCR text:
\begin{itemize}
    \item \textbf{Standard Input:} \texttt{[OCR Text]}
    \item \textbf{Metadata Input:} \texttt{[Dataset:dta19, Language:de, Year:1815] [OCR Text]}
\end{itemize}

We explicitly acknowledge that prepending raw text is a weak conditioning signal for a Seq2Seq architecture. Because we did not register these prefixes as special tokens within the BART tokenizer vocabulary, the model may treat them as arbitrary noise rather than structural metadata. This architectural limitation provides vital context for interpreting BART's quantitative response to metadata.

For \textbf{Qwen}, which relies on instruction tuning, we utilized a system prompt to define the model's task and enforced a strict JSON output schema using the Pydantic framework. This ensures the model explicitly returns a \texttt{corrected\_text} string as output.

% \vspace{0.5em}
\noindent\textbf{Standard Prompt (Baseline):}
\begin{mdframed}[backgroundcolor=gray!10, linewidth=0pt, innertopmargin=10pt, innerbottommargin=10pt, innerleftmargin=10pt]
\small\texttt{Correct the errors in the following OCR text.}\\
\texttt{Return only a JSON object matching this schema: [JSON Schema]}\\
\texttt{OCR Text: [OCR Text]}
\end{mdframed}

\noindent\textbf{Metadata-Enriched Prompt:}
\begin{mdframed}[backgroundcolor=gray!10, linewidth=0pt, innertopmargin=10pt, innerbottommargin=10pt, innerleftmargin=10pt]
\small\texttt{Perform OCR post-correction given this context:\\
Dataset: dta19, Language: de, Year: 1815.}\\
\texttt{Return only a JSON object matching this schema: [JSON Schema]}\\
\texttt{OCR Text: [OCR Text]}
\end{mdframed}
% \vspace{0.5em}

By isolating the metadata variable across these distinct prompting architectures, we aim to quantitatively measure whether external linguistic context and timestamps help models preserve historical orthography.

\subsection{Evaluation Protocol}
Evaluating OCR post-correction models on historical documents requires a precise assessment of how much noise is successfully removed without introducing new generative errors. This study employs standard edit-distance metrics alongside relative improvement scores to quantitatively measure model performance.

\subsubsection{Quantitative Metrics (CER \& WER)}
To measure the absolute volume of OCR noise, we utilize two standard string-matching metrics based on the Levenshtein distance introduced in 1965. In our evaluation pipeline, these metrics are computed using the \texttt{jiwer} Python library:
\begin{itemize}
    \item \textbf{Character Error Rate (CER):} Measures the minimum number of character-level edit operations required to transform the model's output into the curated ground truth. It is formally calculated as:
    \begin{equation}
        CER = \frac{S_{c} + D_{c} + I_{c}}{N_{c}}
    \end{equation}
    where $S_{c}$ is the number of character substitutions, $D_{c}$ is the number of character deletions, $I_{c}$ is the number of character insertions, and $N_{c}$ is the total number of characters in the ground truth reference. CER serves as the primary metric for capturing small orthographic improvements.
    
    \item \textbf{Word Error Rate (WER):} Applies the Levenshtein distance calculation at the token (word) level, measuring the minimum number of word-level edit operations required to align the model's output with the curated ground truth. It is formally calculated as:
    \begin{equation}
        WER = \frac{S_{w} + D_{w} + I_{w}}{N_{w}}
    \end{equation}
    where $S_{w}$ is the number of word substitutions, $D_{w}$ is the number of word deletions, $I_{w}$ is the number of word insertions, and $N_{w}$ is the total number of words in the ground truth reference. WER serves as a secondary metric to evaluate overall readability.
\end{itemize}

\subsubsection{Relative Error Reduction}
While absolute CER and WER provide a snapshot of the final text quality, they do not immediately illustrate the actual value added by the post-correction model relative to the original degraded input. To measure the true efficiency of the models, we calculate the Relative Error Reduction for both characters ($\Delta CER$) and words ($\Delta WER$).

The relative improvement is defined as the percentage reduction in errors from the baseline OCR text to the model's corrected output. For CER, the formula is expressed as:
\begin{equation}
    \Delta CER = \frac{CER_{ocr} - CER_{model}}{CER_{ocr}} \times 100
\end{equation}

Similarly, the Relative Error Reduction for words evaluates the proportional decrease in token-level errors, formulated as:
\begin{equation}
    \Delta WER = \frac{WER_{ocr} - WER_{model}}{WER_{ocr}} \times 100
\end{equation}

A positive percentage means a successful reduction in OCR noise, whereas a negative percentage would indicate that the model introduced more errors than it fixed. By calculating both the absolute scores and the relative reduction percentages, we can systematically compare different model architectures.

\subsubsection{Training and Evaluation Strategies}
To assess the models' domain adaptation and cross-lingual generalization, we use two experimental workflows:

\begin{itemize}
    \item \textbf{Dataset-Level Protocol (Domain-Specific):} In this setup, a model is fine-tuned and evaluated within a single historical dataset. For example, it is trained on the \texttt{dta19\_train} split and evaluated on the \texttt{dta19\_test} split. This provides a baseline for how well the architecture captures a specific historical period, language, and layout.
    
    \item \textbf{Aggregation-Level Protocol (Cross-Domain):} To assess generalized learning, all training datasets are merged into a single multilingual corpus (\texttt{hipe\_aggregated\_train}). A model fine-tuned on this corpus is then evaluated separately on each test set (e.g., \texttt{dta19\_test}, \texttt{icdar2017\_test}, \texttt{french\_test}). This protocol shows whether training on a diverse dataset yields beneficial cross-lingual transfer or instead introduces interference that harms domain-specific accuracy.
\end{itemize}

% Crucially, to prevent data leakage during this aggregated workflow, all dataset-specific test and validation splits were strictly held out prior to aggregation and were never exposed to the model during the training phase.

%========================================== RESULTS =================================================


\section{Results}

\subsection{Quantitative Results}
To evaluate the effectiveness of the different models, we analyzed the Relative Error Reduction for both Character Error Rate ($\Delta CER$) and Word Error Rate ($\Delta WER$) across the aggregated test sets. A positive percentage indicates successful noise reduction, while a negative percentage indicates that the model introduced more errors than it corrected.

As shown in Table \ref{tab:quantitative_results}, the results vary drastically between the highly constrained Seq2Seq architectures and the generative LLMs.

\begin{table}[htbp]
    \centering
    \renewcommand{\arraystretch}{1.1}
    \resizebox{\linewidth}{!}{%
    \begin{tabular}{lllc *{5}{>{\centering\arraybackslash}p{1.8cm}}}
    \toprule
    \multirow{2}{*}{\textbf{Model}} & \multirow{2}{*}{\textbf{Setup}} & \multirow{2}{*}{\textbf{Meta.}} & \multirow{2}{*}{\textbf{Metric}} & \multicolumn{5}{c}{\textbf{Relative Error Reduction ($\Delta\%$)} } \\
    \cmidrule(lr){5-9}
    & & & & \textbf{dta-19} & \textbf{icdar2017} & \textbf{impresso-nzz} & \textbf{\makecell{impresso-\\snippets}} & \textbf{overproof} \\
    \midrule
    % BART-base
    \multirow{4}{*}{\textbf{BART-base}} & \multirow{4}{*}{Finetune} & \multirow{2}{*}{No} & CER & \cR{-4.65} & \cR{-23.64} & \cR{-18.93} & \cR{-9.74} & \cR{-11.72} \\
    & & & WER & \cG{13.46} & \cR{-1.44} & \cG{4.07} & \cG{26.21} & \cG{23.60} \\
    \cmidrule{3-9}
    & & \multirow{2}{*}{Yes} & CER & \cR{-5.04} & \cR{-26.80} & \cR{-24.55} & \cR{-8.15} & \cR{-16.26} \\
    & & & WER & \cG{13.47} & \cR{-2.26} & \cG{2.59} & \cG{26.32} & \cG{23.10} \\
    \midrule
    % BART-large
    \multirow{4}{*}{\textbf{BART-large}} & \multirow{4}{*}{Finetune} & \multirow{2}{*}{No} & CER & \cR{-56.25} & \cR{-60.95} & \cR{-47.21} & \cR{-63.76} & \cR{-24.27} \\
    & & & WER & \cG{5.32} & \cR{-5.82} & \cG{1.07} & \cG{24.34} & \cG{23.90} \\
    \cmidrule{3-9}
    & & \multirow{2}{*}{Yes} & CER & \cR{-92.17} & \cR{-32.23} & \cR{-60.83} & \cR{-69.20} & \cR{-7.21} \\
    & & & WER & \cR{-0.57} & \cR{-0.13} & \cR{-3.80} & \cG{23.76} & \cG{29.69} \\
    \midrule
    % mBART
    \multirow{4}{*}{\textbf{mBART}} & \multirow{4}{*}{Finetune} & \multirow{2}{*}{No} & CER & \cR{-44.34} & \cG{3.44} & \cR{-10.57} & \cR{-29.35} & \cR{-5.73} \\
    & & & WER & \cR{-77.57} & \cG{32.33} & \cG{31.12} & \cR{-66.59} & \cG{7.54} \\
    \cmidrule{3-9}
    & & \multirow{2}{*}{Yes} & CER & \cR{-50.31} & \cR{-8.62} & \cR{-14.89} & \cR{-46.86} & \cR{-21.03} \\
    & & & WER & \cR{-78.88} & \cG{27.59} & \cG{30.77} & \cR{-69.18} & \cG{3.20} \\
    \midrule
    % PyKale
    \multirow{4}{*}{\textbf{PyKale}} & \multirow{4}{*}{-} & \multirow{2}{*}{No} & CER & \cR{-942.52} & \cR{-354.99} & \cR{-689.27} & \cR{-312.45} & \cR{-171.02} \\
    & & & WER & \cR{-219.30} & \cR{-73.38} & \cR{-170.99} & \cR{-85.00} & \cR{-45.65} \\
    \cmidrule{3-9}
    & & \multirow{2}{*}{Yes} & CER & \cR{-1333.27} & \cR{-553.10} & \cR{-969.71} & \cR{-588.26} & \cR{-215.51} \\
    & & & WER & \cR{-278.39} & \cR{-107.33} & \cR{-237.89} & \cR{-118.54} & \cR{-51.97} \\
    \midrule
    % Qwen3-4B Zero-Shot
    \multirow{8}{*}{\textbf{Qwen3-4B}} & \multirow{4}{*}{Zero-Shot} & \multirow{2}{*}{No} & CER & \cR{-177.80} & \cR{-73.89} & \cR{-44.74} & \cR{-104.96} & \cG{33.61} \\
    & & & WER & \cR{-84.94} & \cG{7.86} & \cR{-11.60} & \cR{-53.38} & \cG{27.21} \\
    \cmidrule{3-9}
    & & \multirow{2}{*}{Yes} & CER & \cR{-82.58} & \cG{12.19} & \cR{-7.11} & \cR{-29.01} & \cG{9.88} \\
    & & & WER & \cR{-60.66} & \cG{37.15} & \cR{-0.90} & \cR{-38.71} & \cG{11.28} \\
    \cmidrule{2-9}
    % Qwen3-4B Finetune
    & \multirow{4}{*}{Finetune} & \multirow{2}{*}{No} & CER & \cG{4.40} & \cG{8.35} & \cR{-5.39} & \cG{23.24} & \cG{21.68} \\
    & & & WER & \cG{39.99} & \cG{36.81} & \cG{33.30} & \cB{51.77} & \cB{43.18} \\
    \cmidrule{3-9}
    & & \multirow{2}{*}{Yes} & CER & \cG{8.06} & \cG{14.53} & \cG{1.77} & \cG{23.29} & \cG{23.11} \\
    & & & WER & \cB{40.39} & \cB{40.25} & \cB{41.75} & \cB{50.11} & \cB{42.92} \\
    \midrule
    % Qwen3-8B Zero-Shot
    \multirow{8}{*}{\textbf{Qwen3-8B}} & \multirow{4}{*}{Zero-Shot} & \multirow{2}{*}{No} & CER & \cR{-104.13} & \cG{12.58} & \cR{-4.36} & \cR{-41.73} & \cG{20.56} \\
    & & & WER & \cR{-59.18} & \cB{42.99} & \cG{10.19} & \cR{-34.22} & \cG{23.17} \\
    \cmidrule{3-9}
    & & \multirow{2}{*}{Yes} & CER & \cR{-101.50} & \cG{8.64} & \cR{-5.07} & \cR{-49.59} & \cG{38.54} \\
    & & & WER & \cR{-58.09} & \cB{42.90} & \cG{9.69} & \cR{-39.47} & \cG{29.47} \\
    \cmidrule{2-9}
    % Qwen3-8B Finetune
    & \multirow{4}{*}{Finetune} & \multirow{2}{*}{No} & CER & \cG{0.06} & \cG{26.35} & \cG{16.27} & \cG{32.94} & \cB{46.27} \\
    & & & WER & \cG{39.02} & \cB{44.59} & \cB{40.22} & \cB{48.55} & \cB{55.64} \\
    \cmidrule{3-9}
    & & \multirow{2}{*}{Yes} & CER & \cG{7.74} & \cG{18.06} & \cG{20.07} & \cG{33.37} & \cB{45.66} \\
    & & & WER & \cB{41.44} & \cB{44.52} & \cB{42.77} & \cB{51.68} & \cB{54.65} \\
    \midrule
    % Pleias
    \multirow{4}{*}{\textbf{OCRonos}} & \multirow{4}{*}{Zero-Shot} & \multirow{2}{*}{No} & CER & \cR{-873.99} & \cR{-828.82} & \cR{-36.90} & \cR{-204.92} & \cR{-68.74} \\
    & & & WER & \cR{-158.06} & \cR{-131.61} & \cG{2.19} & \cR{-58.23} & \cR{-6.40} \\
    \cmidrule{3-9}
    & & \multirow{2}{*}{Yes} & CER & \cR{-309.70} & \cR{-154.85} & \cR{-44.05} & \cR{-189.91} & \cR{-9.94} \\
    & & & WER & \cR{-80.39} & \cG{8.03} & \cR{-0.29} & \cR{-50.39} & \cG{11.55} \\
    \bottomrule
    \end{tabular}%
    }
    \caption{Relative Character Error Rate ($\Delta CER$) and Word Error Rate ($\Delta WER$) reductions. Positive values indicate noise reduction; negative values indicate introduced generative errors.}
    \label{tab:quantitative_results}
\end{table}

\noindent\textbf{Key Quantitative Observations}\\
Analyzing the relative error reductions reveals several stark contrasts between model architectures and training paradigms:
\begin{itemize}
    \item \textbf{Seq2Seq CER/WER Divergence:} The BART models showed a striking metric split: they did not improve Character Error Rate, yet \texttt{bart-large} improved Word Error Rate (e.g., $+29.69\%$ on Overproof). This suggests a strong vocabulary mismatch: the model fixes word boundaries and inserts plausible modern words (improving WER) but overwrites archaic spellings with modern subword tokens, sharply worsening CER.
    
    \item \textbf{PyKale Catastrophe:} The PyKale baseline collapsed with CER degrading by over $-900\%$ on some datasets. This originates from its architecture: it was pre-trained with a strict 50-word context window. Applied to our 300-word dynamic chunks, it suffered severe truncation and positional embedding failures, making it effectively incompatible with long-form historical layouts.
    
    \item \textbf{Zero-Shot Degradation:} Zero-shot prompting for historical text correction severely damaged the original OCR, with untuned Qwen3 models and the pre-trained PleIAs/OCRonos model introducing many new errors and large negative scores, indicating prompt engineering alone cannot solve historical OCR repair.
    
    \item \textbf{Fine-Tuning Effectiveness:} QLoRA fine-tuning of Qwen3 produced the strongest results. The fine-tuned \texttt{Qwen3-8B} consistently achieved the largest error reductions, surpassing $50\%$ WER improvement on challenging datasets such as Overproof ($55.64\%$) and Impresso-snippets ($51.68\%$). 
    
    \item \textbf{Metadata Effects:} Historical metadata (year, language, dataset) had minimal effect on Seq2Seq models but substantially improved the smaller generative model (\texttt{Qwen3-4B}) across metrics (see \textbf{Appendix \ref{appendix:charts}, Figure \ref{fig:metadata_qwen4b}}). For the larger \texttt{Qwen3-8B}, metadata effects were inconsistent and sometimes harmful, depending on the dataset.
\end{itemize}

\noindent\textbf{Individual vs. Aggregated Training Protocols}\\
To evaluate transfer learning versus domain interference, we compared our best-performing model (\texttt{Qwen3-8B} with metadata) trained exclusively on specific target datasets (Individual) against training on the entire multilingual corpus (Aggregated). 

As shown in \textbf{Appendix \ref{appendix:charts}, Figure \ref{fig:indiv_vs_agg}}, the aggregated approach is highly competitive and frequently superior. It achieved lower Character Error Rates across almost all datasets, outperforming isolated models on \texttt{dta19} ($3.22\%$ vs. $3.51\%$), \texttt{impresso-nzz} ($4.69\%$ vs. $5.15\%$), \texttt{impresso-snippets} ($2.42\%$ vs. $2.73\%$), and \texttt{overproof-combined} ($5.64\%$ vs. $5.83\%$). 

% MOVE TO APPENDIX ACCORDING TO UNIVERSITY GUIDELINES
% \begin{figure}[H]
%     \centering
%     \includegraphics[width=0.95\textwidth]{figures/metadata_impact_qwen4b.png}
%     \caption{Impact of metadata integration on the fine-tuned \texttt{Qwen3-4B} model. Supplying explicit historical context (green bars) consistently gives lower absolute Character Error Rates and Word Error Rates compared to the standard prompt without metadata (blue bars).}
%     \label{fig:metadata_qwen4b}
% \end{figure}

Despite minor word-level regressions (e.g., a slight WER increase on \texttt{impresso-snippets} from $9.79\%$ to $10.90\%$), the strong cross-lingual and cross-domain transfer learning validates using a single, unified generative model rather than a fragmented ensemble of dataset-specific models.

\noindent\textbf{A Note on Statistical Variance:} 
It is important to acknowledge that due to the heavy computational constraints of fine-tuning large language models on a single GPU, these quantitative results represent single training runs. While the relative percentage improvements are substantial enough to establish clear architectural trends, future work with an expanded timeframe should include multi-run averaging (e.g., 5-fold cross-validation) to establish strict confidence intervals and measure variance.


\subsection{Qualitative Results and Error Taxonomy}
While quantitative metrics (CER and WER) capture the absolute volume of edit operations, they treat all character changes equally. They cannot distinguish between a model that modernizes an archaic word (a loss of historical fidelity) and a model that hallucinates a completely new phrase. To address this, we categorize the model outputs into a unified qualitative error taxonomy (Table \ref{tab:unified_qualitative}).

\begin{table}[H]
    \centering
    \renewcommand{\arraystretch}{1.5}
    \resizebox{\linewidth}{!}{%
    \begin{tabular}{>{\raggedright\arraybackslash}p{2.8cm} >{\raggedright\arraybackslash}p{4cm} >{\raggedright\arraybackslash}p{4.5cm} >{\raggedright\arraybackslash}p{4cm}}
    \toprule
    \textbf{Error Taxonomy} & \textbf{Original OCR} & \textbf{Model Output} & \textbf{Ground Truth} \\
    \midrule
    \multicolumn{4}{c}{\textbf{Sequence-to-Sequence Models (BART)}} \\
    \midrule
    
    \textbf{1. Stylistic Rewrite} & 
    It has \textbf{fillen} under our notice... \newline \textbf{moëlle, poëte} & 
    It has \textbf{\textcolor{red}{come}} under our notice... \newline \textcolor{red}{\textbf{moelle, poète}} \newline \textit{(Model: BART-base)} & 
    It has \textbf{\textcolor{bestgreen}{fallen}} under our notice... \newline \textcolor{bestgreen}{\textbf{moëlle, poëte}} \\
    
    \textbf{2. Statistical Guessing} & 
    THE LATE \textbf{Mn.} HERER & 
    THE LATE \textbf{\textcolor{red}{MRS.}} HERER \newline \textit{(Model: BART-base)} & 
    THE LATE \textbf{\textcolor{bestgreen}{MR.}} HEBER \\
    
    \textbf{3. Semantic Surrender} & 
    nor to \textbf{expiscate} an invitation... & 
    nor to \textbf{\textcolor{red}{explicate}} an invitation \textbf{\textcolor{red}{in any hint}}... \newline \textit{(Model: BART-large)} & 
    nor to \textbf{\textcolor{bestgreen}{expiscate}} an invitation \textbf{\textcolor{bestgreen}{from him}}... \\
    
    \midrule
    \multicolumn{4}{c}{\textbf{Generative LLMs (Qwen3, PleIAs/OCRonos)}} \\
    \midrule
    
    \textbf{4. Contextual Hallucination} & 
    \textbf{Itara.}-The Golconda... \newline
    22 octobre \textbf{in séré} ce matin... & 
    \textbf{\textcolor{red}{Hogshead}} - The Golconda... \newline
    22 octobre \textbf{\textcolor{red}{hier}} ce matin... & 
    \textbf{\textcolor{bestgreen}{Rum.}} - The Golconda... \newline
    22 octobre \textbf{\textcolor{bestgreen}{inséré}} ce matin... \\
    
    \textbf{5. Semantic Override} & 
    supplies from \textbf{Tunk} Ceylon... \newline
    dans le Sahel des \textbf{Zlass}... & 
    supplies from \textbf{\textcolor{red}{Burma}} Ceylon... \newline
    dans le Sahel des \textbf{\textcolor{red}{Kassis}}... & 
    supplies from \textbf{\textcolor{bestgreen}{Junk}} Ceylon... \newline
    dans le Sahel des \textbf{\textcolor{bestgreen}{Zlass}}... \\
    
    \textbf{6. Aggressive Modernization} & 
    at \textbf{5} dollars to \textbf{74} dollars... \newline
    délégués et \textbf{suppléans}... & 
    at \textbf{\textcolor{red}{\$5}} dollars to \textbf{\textcolor{red}{\$74}} dollars... \newline
    délégués et \textbf{\textcolor{red}{suppléants}}... & 
    at \textbf{\textcolor{bestgreen}{5}} dollars to \textbf{\textcolor{bestgreen}{74}} dollars... \newline
    délégués et \textbf{\textcolor{bestgreen}{suppléans}}... \\
    
    \midrule
    \multicolumn{4}{c}{\textbf{Impact of Metadata Conditioning (Qwen3 Models)}} \\
    \midrule

    \textbf{7. Prevents Cross-Lingual Translation} & 
    \textbf{unter} Hadrian... & 
    \textit{No Meta:} \textbf{\textcolor{red}{The Emperor}} Hadrian... \newline
    \textit{With Meta:} \textbf{\textcolor{bestgreen}{unter}} Hadrian... & 
    \textbf{\textcolor{bestgreen}{unter}} Hadrian... \\

    \textbf{8. Prevents Bigram Traps} & 
    \textbf{pnfu} ano Portugal. The advices from Spain... & 
    \textit{No Meta:} \textbf{\textcolor{red}{France}} and Portugal... \newline
    \textit{With Meta:} \textbf{\textcolor{bestgreen}{Spain}} and Portugal... & 
    \textbf{\textcolor{bestgreen}{Spain}} and Portugal... \\

    \textbf{9. Prevents Generic Substitution} & 
    aus dem \textbf{Ober land} berichtet wird... & 
    \textit{No Meta:} aus dem \textbf{\textcolor{red}{Osten}} berichtet... \newline
    \textit{With Meta:} aus dem \textbf{\textcolor{bestgreen}{Ober¬land}} berichtet... & 
    aus dem \textbf{\textcolor{bestgreen}{Ober¬land}} berichtet wird... \\
    
    \bottomrule
    \end{tabular}%
    }
    \caption{Unified Qualitative Error Taxonomy. Seq2Seq models (top) exhibit stylistic paraphrasing and statistical guessing, fundamentally acting as modernizing translators. Generative LLMs (middle) prioritize semantic context over visual reality. The bottom section demonstrates how explicit metadata conditioning helps generative models prevent cross-lingual translation and statistical bigram traps.}
    \label{tab:unified_qualitative}
\end{table}

\noindent\textbf{Seq2Seq Models: Lexical Hallucinations}\\
Although the Seq2Seq models (BART) are constrained as denoising auto-encoders, they showed a severe bias toward their modern pre-training distributions. Instead of restoring texts in a faithful manner, they frequently act as modernizing translators. As highlighted in the taxonomy, this bias is illustrated through stylistic rewrites, statistical guesses despite clear visual clues, and completely surrendering to modern synonyms when encountering rare historical vocabulary.

\noindent\textbf{Generative Models: Hallucination and Over-Correction}\\
The inference results reveal that errors in the generative model occur when internal semantic probabilities outweigh strict visual evidence. Instead of repairing individual characters, LLMs commonly impose geographic priors, hallucinate contextually plausible but erroneous historical entities, and aggressively force modern formatting (such as modern decimal systems or updated pluralizations) upon the text. 

\noindent\textbf{Impact of Metadata Conditioning}\\
Explicit document metadata acts as an anchor against the unconstrained generative biases of LLMs. Without metadata, the model falls into classic statistical association traps, defaulting to highly frequent modern token sequences (converting the corrupted string "\texttt{pnfu ano Portugal}" into the common bigram "\texttt{France and Portugal}"). Similarly, unconditioned models default to generic substitutions, swapping localized geographical text ("\texttt{Oberland}") for common directional tokens ("\texttt{Osten}"). Most severely, lacking language constraints can trigger complete cross-lingual hallucinations, where the model attempts to translate the text into English (hallucinating "The Emperor" for the German word "unter"). When metadata conditioning is introduced, the model successfully exploits temporal, language, and collection signatures to anchor its predictions directly to the historical visual evidence.
\subsection{Best Configuration Refinement}

\noindent\textbf{Motivation and Setup}\\
After establishing that the fine-tuned \texttt{Qwen3-8B} with metadata was our best-performing model, we conducted a targeted ablation study to push its performance even further. We tested three specific changes on the aggregated training protocol to see which parameters give the best results:

\begin{itemize}
    \item \textbf{More Epochs:} Increasing the training duration from 5 to 10 epochs.
    \item \textbf{Better Prompt:} Updating the system prompt with strict, explicit instructions and concrete examples of historical OCR errors to prevent them from modernizing our texts. The full content of the prompt can be found in \textbf{Appendix \ref{appendix:prompt}}.
    \item \textbf{Adjust Hyperparameters:} Replace the standard learning rate scheduler with a \texttt{cosine} decay scheduler and significantly increase the LoRA scaling factor ($\alpha=128$ instead of $16$) to force the model to adopt the fine-tuning weights more aggressively.
\end{itemize}

\noindent\textbf{Ablation Results}\\
We tested these modifications individually and then combined the best approaches. The absolute Character Error Rate (CER) and Word Error Rate (WER) are presented in Table \ref{tab:ablation_results}.

\begin{table}[htbp]
    \centering
    \renewcommand{\arraystretch}{1.2}
    \resizebox{\linewidth}{!}{%
    \begin{tabular}{l c c c c c}
    \toprule
    \textbf{Configuration} & \textbf{dta-19} & \textbf{icdar2017} & \textbf{impresso-nzz} & \textbf{impresso-snip} & \textbf{overproof} \\
    & CER / WER & CER / WER & CER / WER & CER / WER & CER / WER \\
    \midrule
    Initial (Baseline) & 3.22 / 10.94 & 2.49 / 7.37 & 4.69 / 13.19 & 2.42 / 10.90 & 5.64 / 13.79 \\
    \midrule
    10 Epochs & \cR{3.81} / \cR{11.55} & \cR{2.88} / \cR{8.46} & \cR{5.48} / \cR{13.93} & \cR{3.58} / \cR{12.06} & \cR{5.81} / \cR{14.07} \\
    Better Prompt & \cG{2.95} / \cG{10.17} & \cB{2.11} / \cB{7.15} & \cG{4.36} / \cG{12.15} & \cB{2.22} / \cR{10.94} & \cR{6.52} / \cR{14.95} \\
    Hyperparams ($\alpha=128$) & \cG{3.15} / \cG{10.22} & \cR{2.59} / \cR{7.63} & \cG{4.49} / \cG{13.09} & \cR{2.44} / \cB{10.18} & \cR{8.25} / \cR{16.41} \\
    Combined & \cB{2.87} / \cB{9.69} & \cG{2.45} / \cG{7.62} & \cB{3.98} / \cB{11.56} & \cR{2.44} / \cG{10.71} & \cB{5.38} / \cR{14.54} \\
    \bottomrule
    \end{tabular}%
    }
    \caption{Ablation results on the \texttt{Qwen3-8B} (with metadata) aggregated model. Green indicates improvement over the baseline, dark green indicates the best result, and red indicates degradation.}
    \label{tab:ablation_results}
\end{table}

\noindent\textbf{Key Observations}\\
This ablation study reveals exactly which configurations offer the highest performance gains for OCR post-correction:

\begin{itemize}
    \item The 10-epoch run degraded performance across almost all datasets. This indicates that the model was overfitting, memorizing the training noise rather than generalizing.
    \item The highly detailed prompt alone improved CER on 4 out of 5 datasets. By explicitly providing examples of common OCR errors and commanding the model to act as an archivist, we prevent semantic hallucinations. This shows that prompt engineering is a mandatory optimization tool.
    \item By increasing the LoRA $\alpha$ to 128, it forces the model to learn the training data much more aggressively. While this helped on \texttt{impresso-snippets}, it caused catastrophic degradation on \texttt{overproof} (CER jumping from 5.64 to 8.25). We hypothesize that \texttt{overproof}'s English text is already well-represented in Qwen3's pre-training, so the aggressive weight updates and extended cosine warm-up caused unnecessary semantic drift.
    \item Combining the better prompt with the modified hyperparameters gave the best overall results. For example, the absolute WER on \texttt{impresso-nzz} decreased from 13.09\% to 11.56\%. However, on \texttt{icdar2017}, the prompt alone was slightly better, suggesting that future configurations should utilize the advanced prompt while keeping the LoRA $\alpha$ at a more conservative baseline.
\end{itemize}

\subsection{Competition Leaderboard}

As the ultimate external validation of our proposed methodology, our \texttt{Qwen3-8B} refined combined configuration was submitted to the official ICDAR 2026 HIPE-OCRepair scorer. The official competition metric is the character-level Minimum Edit Rate (cMER), a normalized edit distance computed across all test set outputs. Table \ref{tab:leaderboard} displays the evaluation leaderboard across all 8 hidden test sets.

\begin{table}[H]
    \centering
    \renewcommand{\arraystretch}{1.3}
    \resizebox{\linewidth}{!}{%
    \begin{tabular}{c l c c c c}
    \toprule
    \textbf{Rank} & \textbf{System} & \makecell{\textbf{Overall} \\ \textbf{cMER $\downarrow$}} & \textbf{95\% CI} & \makecell{\textbf{Overall Pref} \\ \textbf{Macro $\uparrow$}} & \textbf{95\% CI} \\
    \midrule
    1 & \texttt{bnf-mistral\_hipe-ocrepair-bench\_v0.9\_run1} & 0.0050 & [0.004, 0.007] & 0.9000 & [0.835, 0.952] \\
    2 & \texttt{bnf-mistral\_hipe-ocrepair-bench\_v0.9\_run3} & 0.0071 & [0.005, 0.009] & 0.8737 & [0.803, 0.931] \\
    3 & \texttt{bnf-mistral\_hipe-ocrepair-bench\_v0.9\_run2} & 0.0090 & [0.007, 0.011] & 0.8739 & [0.791, 0.942] \\
    4 & \texttt{blocr\_hipe-ocrepair-bench\_v0.9\_run1} & 0.0106 & [0.008, 0.013] & 0.7028 & [0.574, 0.816] \\
    \rowcolor{green!10}
    \textbf{5} & \textbf{\texttt{TRANMINHDUONG\_hipe-ocrepair-bench\_v0.9\_run1}} & \textbf{0.0167} & \textbf{[0.013, 0.021]} & \textbf{0.5694} & \textbf{[0.398, 0.732]} \\

    6 & \texttt{l3i\_hipe-ocrepair-bench\_v0.9\_run1} & 0.0176 & [0.014, 0.022] & 0.3591 & [0.229, 0.486] \\
    \midrule
    9 & \texttt{baseline-no-correction\_hipe-ocrepair} & 0.0226 & [0.019, 0.026] & 0.0000 & [0.000, 0.000] \\
    \bottomrule
    \end{tabular}%
    }
    \caption{Leaderboard for the ICDAR 2026 HIPE-OCRepair competition. Our submission \texttt{TRANMINHDUONG} is highlighted, successfully outperforming the official no-correction baseline and that of \texttt{l3i}.}
    \label{tab:leaderboard}
\end{table}

Our submission (\texttt{TRANMINHDUONG}) outperformed the official \texttt{baseline-no-correction} system (Rank 9) and the \texttt{l3i} submission (Rank 6), achieving a competitive 5th-place overall ranking. The official competition report highlights the substantial resources and complex pipelines used by the higher-ranking teams, which helps put our results into perspective. The top three runs by \texttt{BNF-MISTRAL}, for example, relied on a much larger 24-billion-parameter model (Mistral Small 3) that was further pre-trained on historical corpora and fine-tuned on synthetic data. Their method also used a resource-intensive, multi-step inference pipeline with iterative judge-and-retry loops and conditional routing to even larger models. The 4th-place \texttt{BLOCR} system followed a model-switching strategy, using DeepSeek for English and Gemma 3 for French and German. Even the 6th-place \texttt{l3i} team relied on a larger 14-billion-parameter model (Qwen3-14B) and, notably, did not include document metadata in their prompts.

Against this backdrop, our results are highly encouraging. They provide external validation for our core hypothesis: a single, lightweight 8-billion-parameter model, adapted via Parameter-Efficient Fine-Tuning (QLoRA) on a single GPU. Guided by explicit metadata constraints, it can deliver robust OCR post-correction that competes with far more complex and computationally expensive systems in an international competition setting.

%========================================== DISCUSSION ===============================================
\section{Discussion}

\subsection{Trade-off Analysis: Correction vs. Preservation}
The main challenge of historical OCR post-correction is that it is not simply a denoising problem. A correction model must decide what to change and what to preserve. Our results reveal a contradiction between developing reading fluency and maintaining archival authenticity. Although a model can successfully reduce Word Error Rate by injecting modern vocabulary, it may introduce character-level modernization that damages historical fidelity. In this domain, fluency is not synonymous with fidelity. This trade-off dictates that models cannot be optimized for standard language modeling objectives without destroying the unique historical character of the source texts.

\subsection{Comparison Between Models}
The differences in behavior between Seq2Seq architectures and generative LLMs are the most striking findings of this study. Despite being architecturally limited to the input sequence, Seq2Seq models (BART) demonstrated a strong bias toward their modern pre-training distributions.  

On the other hand, decoder-only LLMs (Qwen3) proved more flexible and achieved stronger results after fine-tuning. However, this generative flexibility inherently came at the cost of semantic drift and hallucination. While BART fails by stubbornly updating the historical text, unconstrained LLMs fail by intentionally inventing historically plausible but visually unsupported content. This indicates that zero-shot prompting is fundamentally insufficient for this task and generative models require strict, preservation-oriented fine-tuning to connect their semantic knowledge to the real meaning behind the OCR string.

\subsection{Limitations}
This study has several limitations that could provide directions for future research. Firstly, due to the significant computational resources required to fine-tune large language models on a single GPU, the quantitative results represent single training runs rather than multi-run averages. Secondly, in our experiments, document metadata was inserted merely as a textual prefix or prompt inclusion. A more structured metadata encoding strategy might give stronger and more consistent conditioning signals for Seq2Seq models. Finally, the models operated purely on text input. They did not have access to images of the document on which OCR was performed, meaning they could not verify their corrections against the original scanned document crops.

%========================================== CONCLUSION =================================================
\section{Conclusion}

This internship report investigated whether parameter-efficient fine-tuning on domain-specific data allows modern LLMs to reduce OCR noise while preserving historical spellings better than traditional Seq2Seq models. Our quantitative and qualitative results confirm the hypothesis, but with critical limitations.

Our results demonstrate that QLoRA fine-tuned generative models (Qwen3-4B and Qwen3-8B) significantly outperform standard encoder-decoder models (BART) across multiple multilingual historical datasets of the HIPE-OCRepair 2026 benchmark, achieving the highest relative error reductions. However, we also found that generative LLMs inherently suffer from semantic drift, prioritizing contextual plausibility over visual fidelity. To successfully navigate the trade-off between text correction and historical preservation, LLMs require a combination of aggregated multilingual fine-tuning, explicit metadata conditioning, and robust prompt engineering. 

Ultimately, this study proves that historical OCR post-correction is a unique preservation task. Future work must bridge the gap between text-based generation and visual reality by exploring multimodal architectures that ground LLM corrections directly in the historical document images.


\subsection*{Personal Reflection}
This internship has been an truly enriching experience that went far beyond the scope of OCR post-correction. Working within the dynamic environment of the L3i laboratory gave my first genuine impression of an active research community. I want to express my deepest gratitude to my supervisors, Dr. Emanuela Boros and Prof. Antoine Doucet, whose continuous mentorship and patience guided me through the entire lifecycle of a research study. From reading countless papers for literature review, designing methodologies, implementing training pipelines, and writing results in a formal academic format, their support was instrumental in teaching me how to navigate the research world.

On a technical level, I gained hands-on experience with large language model fine-tuning, prompt engineering, and multilingual NLP evaluation, skills which I had previously only encountered in theoretical lectures. Perhaps the most valuable technical lesson was learning how to bridge that gap between understanding a theoretical concept and actually implementing it under real-world hardware constraints. Solving these practical bottlenecks has equipped me with the necessary foundation as I prepare to transition into my upcoming role as an AI Engineer when I return to Vietnam.

The most challenging aspect of the research was learning how to interpret negative results. At first, I found it difficult to present the degraded performance of models like BART-large or PyKale without feeling that my methodology was inaccurate. However, through insightful discussions with Dr. Boros, I came to the realization that a well-explained negative result is just as valuable as a positive one.

At the end of the day, my time at the L3i laboratory did not just make me a better programmer; it fundamentally improved my problem-solving mindset, bridging the gap between my academic studies and my future career in AI engineering.

% ==========================================
% BIBLIOGRAPHY & APPENDIX
% ==========================================
\newpage
\titleformat{name=\section,numberless}[display]
  {\normalfont\huge\bfseries\filcenter}
  {}
  {0pt}
  {} 
  []

\bibliographystyle{unsrturl}
\bibliography{references}
\addcontentsline{toc}{section}{References}

\appendix

\section{Performance Charts}
\label{appendix:charts}

\begin{figure}[H]
    \centering
    \includegraphics[width=0.95\textwidth]{figures/metadata_impact_qwen4b.png}
    \caption{Impact of metadata integration on the fine-tuned \texttt{Qwen3-4B} model. Supplying explicit historical context (green bars) consistently gives lower absolute Character Error Rates and Word Error Rates compared to the standard prompt without metadata (blue bars).}
    \label{fig:metadata_qwen4b}
\end{figure}

\begin{figure}[H]
    \centering
    \includegraphics[width=0.95\textwidth]{figures/protocol_comparison.png}
    \caption{Absolute CER and WER performance of the fine-tuned \texttt{Qwen3-8B} (with metadata). The chart compares the initial baseline noise against models trained individually on specific datasets versus a single model trained on the aggregated corpus.}
    \label{fig:indiv_vs_agg}
\end{figure}

\section{Detailed Preservation-Oriented Prompt}
\label{appendix:prompt}

Below is the advanced system prompt utilized during the Best Configuration Refinement experiment to enforce historical fidelity and prevent generative modernization.

\begin{mdframed}[backgroundcolor=gray!10, linewidth=0pt, innertopmargin=10pt, innerbottommargin=10pt, innerleftmargin=10pt, innerrightmargin=10pt]
\small\ttfamily\raggedright
You are an expert OCR post-correction system for historical newspapers in many languages. \\
Your task is to correct the OCR hypothesis while preserving the original historical wording, spelling, punctuation style, paragraph flow, and meaning as much as possible. \\
You are given the document metadata and the OCR hypothesis. You must output only the corrected transcription unit formatted as JSON.\\

\noindent Important principles:\\
- Correct OCR errors caused by character confusion, broken words, incorrect punctuation, spacing errors, and misrecognized letters.\\
- Preserve historical spelling and period-specific language. Do not modernize vocabulary. Do not rewrite sentences for fluency. Do not paraphrase.\\
- Do not add information that is not supported by the OCR text. Do not remove content unless it is clearly an OCR artefact.\\
- Keep the text as one continuous transcription unit unless the input clearly contains paragraph breaks.\\
- Preserve names, initials, titles, abbreviations carefully, and capitalization when historically plausible.\\
- Be especially careful with 19th-century punctuation, long sentences, hyphenated compounds, and editorial address forms such as "Mr. Editor" or "Sir".\\

\noindent Document context:\\
- Dataset: \{dataset\}, Language: \{lang\}, Year: \{year\_str\}\\

\noindent Expected correction behavior for this document:\\
The OCR hypothesis comes from a historical newspaper passage in \{lang\}. It contains many typical OCR errors:\\
- Misread letters: "di position" $\rightarrow$ "disposition", "phrenotogized" $\rightarrow$ "phrenologized"\\
- Confused punctuation: "R '" $\rightarrow$ "R.", "prove '" $\rightarrow$ "prove ?"\\
- Misread words: "nliat docs" $\rightarrow$ "what does", "site" $\rightarrow$ "she", "tire" $\rightarrow$ "the"\\
- Broken or merged words: "thesameroutineof" $\rightarrow$ "the same routine of"\\
- Wrong substitutions caused by old print: "lie" may actually be "be", "tlie" may be "the", "tlmt" may be "that"\\
- Hyphenation and line-break artefacts: repair only when clearly caused by OCR\\

\noindent Correction strategy:\\
1. Read the full OCR hypothesis before correcting.\\
2. Infer intended words from local context and historical newspaper style.\\
3. Correct high-confidence OCR errors.\\
4. Leave uncertain historical or rare words unchanged unless the OCR form is clearly impossible.\\
5. Preserve the author's rhetorical style, including irony, long sentences, and formal address.\\
6. Do not hallucinate missing passages.\\

\noindent Return ONLY a JSON object matching this schema: \{schema\_str\}
\end{mdframed}

\end{document}
